% =====================================================================
\section{Beyond the Paper: Robustness, Sensitivity and Findings}
\label{sec:beyond}
% =====================================================================

The paper shows two single runs. To know whether its conclusions hold in general, and where they stop holding, we ran
three Monte-Carlo studies (300 runs each, the same random draws for all three controllers, fixed seed) and a
sensitivity study on the parameters the paper does not state.

% ---------------------------------------------------------------------
\subsection{Monte-Carlo studies}
\label{sec:mc}
% ---------------------------------------------------------------------

\begin{center}\small
\begin{tabularx}{\linewidth}{@{}l X@{}}
\toprule
\textbf{Study} & \textbf{Random draws (per run)}\\
\midrule
A --- paper regime & attitude uniform inside 90\,\% of the envelope ($|q_i(0)|\le0.9\rho_{q0}$), $|\omega_i(0)|\le0.05$\,rad/s, paper disturbance\\
B --- disturbance amplitude & as A, disturbance multiplied by a factor drawn from $U[0.5,\,3]$\\
C --- initial-rate stress & as A, but $|\omega_i(0)|$ up to $0.9\rho_{\omega0}=0.36$\,rad/s\\
\bottomrule
\end{tabularx}
\end{center}

\resfig{fig11_monte_carlo.png}{\linewidth}{Fraction of violation-free runs. A: paper regime with 95\,\% Wilson confidence
intervals. B: by disturbance scale. C: by initial rate.}{fig:mc}

\begin{table}[H]
\centering\small
\setlength{\tabcolsep}{3.5pt}
\caption{Monte-Carlo results ($N=300$ per controller and study). ``Violation-free'' = never outside either envelope.}
\label{tab:mc}
\begin{tabular}{@{}llcccc@{}}
\toprule
\textbf{Study} & \textbf{Controller} & \textbf{violation-free} & \textbf{95\,\% CI} & \textbf{median settling} & \textbf{mean torque energy}\\
\midrule
A & constrained (paper) & \textbf{97.0\,\%} & [94.4, 98.4] & 4.85\,s & 27.9\\
  & constrained (symmetric) & 97.7\,\% & [95.3, 98.9] & 4.89\,s & 28.5\\
  & conventional PD & 0.0\,\% & [0.0, 1.3] & 5.93\,s & 36.0\\
\midrule
B & constrained (paper) & 21.3\,\% & [17.1, 26.3] & 5.05\,s & 29.1\\
  & constrained (symmetric) & \textbf{97.3\,\%} & [94.8, 98.6] & 5.03\,s & 29.8\\
  & conventional PD & 3.0\,\% & [1.6, 5.6] & 6.00\,s & 37.2\\
\midrule
C & constrained (paper) & 50.3\,\% & [44.7, 56.0] & 6.72\,s & 76.5\\
  & constrained (symmetric) & 51.3\,\% & [45.7, 56.9] & 6.68\,s & 89.6\\
  & conventional PD & 0.0\,\% & [0.0, 1.3] & 7.65\,s & 59.3\\
\bottomrule
\end{tabular}
\end{table}

\begin{projbox}[what the three studies say]
\begin{enumerate}[leftmargin=*]
  \item \textbf{A confirms the paper in general, not just for its two initial states.} In the paper's operating regime
        97\,\% of all random starts stay inside both envelopes with the paper's law; conventional PD fails in every run
        (it leaves the rate envelope both during the transient and after it has closed to $\rho_{\omega T}=0.005$\,rad/s,
        where the disturbance keeps its residual rate too large). The
        constrained law also settles about 1\,s earlier and uses 22\,\% less torque energy on average.
  \item \textbf{B: the paper's one-sided barrier is sensitive to the disturbance size; the symmetric barrier is not.} Up to
        the paper's disturbance amplitude every form works (100\,\% for both constrained forms in the $0.5$--$1\times$ bin).
        Above it, the paper form drops to 23\,\% ($1$--$1.5\times$) and 0\,\% beyond $1.5\times$, while the symmetric form stays at
        96--97\,\% up to $3\times$. In the runs above $1.5\times$, \emph{every} sample outside the envelope is at the \emph{lower} wall of the rate
        envelope ($\omega_i<-\rho_\omega$), and all of them occur late ($t>11$\,s, when the tube is narrow), where the paper's barrier
        gain stays near $\beta/4$: exactly the one-sided property of \cref{sec:barrierform}.
  \item \textbf{C: no controller can beat physics.} With $J\approx17$\,kg\,m$^2$ and 2.5\,N\,m the body can decelerate by only
        about $0.15$\,rad/s$^2$ (Mathematical Models, block 1). An initial rate of 0.3\,rad/s needs about 2\,s to stop, while
        the rate envelope is already closing. Both constrained forms degrade identically with the initial rate (100\,\%
        below 0.12\,rad/s, about 35\,\% at 0.33\,rad/s): this is a feasibility limit of the envelope, not a controller defect.
\end{enumerate}
\end{projbox}

% ---------------------------------------------------------------------
\subsection{Sensitivity to the parameters the paper does not give}
\label{sec:sens}
% ---------------------------------------------------------------------

\begin{table}[H]
\centering\small
\caption{Constrained law (paper form), Part-2 initial state, different $T_s$ and $\rho_{q0}$.}
\label{tab:sens}
\begin{tabular}{@{}cccccc@{}}
\toprule
$T_s$ (s) & $\rho_{q0}$ & outside $\rho_q$ & outside $\rho_\omega$ & torque at limit & $\max|\varepsilon_1|$\\
\midrule
15 & 0.36 (ours) & 0.0\,\% & 0.0\,\% & 6.3\,\% & 0.868\\
15 & 0.50 & 0.0\,\% & 0.0\,\% & 7.7\,\% & 0.625\\
10 & 0.50 & 0.0\,\% & 9.6\,\% & 7.4\,\% & 0.813\\
10 & 0.36 & 65.9\,\% & 55.8\,\% & 71.3\,\% & 6.554\\
\bottomrule
\end{tabular}
\end{table}

\resfig{fig10_infeasible_envelope.png}{0.95\linewidth}{An infeasible envelope ($T_s=10$\,s with the tight $\rho_{q0}=0.36$, Part-2
state): the barrier asks for more torque than 2.5\,N\,m, the torque stays pinned at the limit and the constraint is
lost.}{fig:infeasible}

\begin{warn}[Lesson: the envelope must be feasible for the actuator]
The paper presents $\rho_0$, $\rho_T$, $T_s$ as free design choices ``based on the requirements''. They are not
independent of $\umax$: if the tube closes faster than the saturated actuator can turn the body, no gain can keep the state
inside, and the barrier only drives the actuator into permanent saturation. This is the same limit as Monte-Carlo C, and it
is the main design rule we took to the X3, whose inertia and torque authority are completely different.
\end{warn}

% ---------------------------------------------------------------------
\subsection{Findings: what the paper does not say}
\label{sec:findings}
% ---------------------------------------------------------------------

\begin{center}\small
\begin{longtable}{@{}l p{6.4cm} p{6.6cm}@{}}
\toprule
\textbf{\#} & \textbf{Finding} & \textbf{How we handled it}\\
\midrule\endhead
P1 & Four values are not given: $c$, $\varsigma$, the Part-2 settling time and a usable Part-2 $\rho_{q0}$. & $c=0.5$ (satisfies Theorem 1, $c<0.964$), $\varsigma=0.02$, $T_s=15$\,s, $\rho_{q0}=0.36$; sensitivity measured (\cref{tab:sens}).\\
P2 & Part 2 starts on its own boundary ($|q_1(0)|=|q_3(0)|=0.3=\rho_{q0}$), so $(1-\varepsilon)^{-2}=\infty$ and Remark~7 is violated. & $\rho_{q0}=1.2\max|q_i(0)|=0.36$.\\
P3 & The barrier $(1-\varepsilon)^{-2}$ is one-sided. & Paper form is the reference; the symmetric form is reported as a variant; Monte-Carlo B shows the difference.\\
P4 & Envelope parameters must be feasible for $\umax$ and $J$. & Sensitivity study and Monte-Carlo C; used to size the X3 envelope.\\
P5 & Theorem 1 proves convergence to the origin, not that $|\varepsilon|<1$ for all $t$. & Constraint satisfaction is reported as measured, never as guaranteed.\\
P6 & Eq.~(15) is printed with $(1-q_0^2)$ and $\ttil^2/\eta_1$; the derivative (17) follows from $(1-q_0)^2$ and $\tfrac{\eta_1}{2}\ttil^2$. & Written in the consistent form (Mathematical Models, Theorem 1).\\
P7 & With the boundary layer (22) and a leakage that vanishes, $\that$ slowly drifts upward. & Observed (\cref{fig:theta}); on the X3 the leakage has a floor and $\that$ an upper bound.\\
\bottomrule
\end{longtable}
\end{center}
